\documentclass[11pt]{article}
\usepackage{mathblatt}
% fehlender Baustein (nicht in der Anleitung): Quadrat mit Rand für Sachaufgaben (Nr. 44)
\newcommand{\quadratrand}[2]{\begin{tikzpicture}\draw[thick] (0,0) rectangle (3.2,3.2);\fill[gray!15] (0.6,0.6) rectangle (2.6,2.6);\draw (0.6,0.6) rectangle (2.6,2.6);\node at (1.6,1.6) {Foto};\draw[<->] (0.6,0.3) -- (2.6,0.3) node[midway,below] {$#1$};\draw[<->] (0,1.6) -- (0.6,1.6); \node[left] at (0,1.6) {#2};\end{tikzpicture}}
\def\mitzone{}
\begin{document}
\blattkopf{Quadratische Gleichungen}{Gesamt}
\weit
\verzeichniszeile{\verz{zone}{Kennst du schon (Nr. 1–11)} \verztrenn \verz{e1}{1 Wurzelziehen und Lösbarkeit (Nr. 12–23)} \verztrenn \verz{e2}{2 Normalform und p-q-Formel (Nr. 24–39)} \verztrenn \verz{e3}{3 Sachaufgaben (Nr. 40–46)} \verztrenn \verz{e4}{Ausblick 4 Satz vom Nullprodukt (Nr. 47–55)} \verztrenn \verz{abhaken}{Das kann ich}}
% Zone „Kennst du schon" – Hauptnummern 1 bis 11
\einheitenkopf[zone][Kennst du schon]{Das kennst du schon}
\setcounter{aufgabe}{0}

\begin{aufgabe}{Ich kann Zahlen quadrieren, auch negative Zahlen und Dezimalzahlen.}
Berechne.
\begin{teilezwei}
\tz $6^2 = {}$ \leerfeld & \tz $10^2 = {}$ \leerfeld \\
\tz $(-8)^2 = {}$ \leerfeld & \tz $2{,}5^2 = {}$ \leerfeld \\
\tz $(-0{,}5)^2 = {}$ \leerfeld & \\
\end{teilezwei}
\end{aufgabe}

\begin{aufgabe}{Ich kann mit Klammern und negativen Zahlen rechnen: Klammer zuerst, Punkt vor Strich.}
Berechne.
\begin{teilezwei}
\tz $3 \cdot (2 + 4) = {}$ \leerfeld & \tz $20 - 4 \cdot 3 = {}$ \leerfeld \\
\tz $(-2) \cdot 6 = {}$ \leerfeld & \tz $(-4) \cdot (-4 + 9) = {}$ \leerfeld \\
\tz $2 \cdot (-3)^2 - 10 = {}$ \leerfeld & \\
\end{teilezwei}
\end{aufgabe}

\begin{aufgabe}{Ich kann eine lineare Gleichung lösen.}
Löse die Gleichung.
\begin{gleichungsraster}[2]
\gl{x + 4 = 10} & \gl{3x = 12} \\
\gl{x - 2{,}5 = 4} & \gl{x + 6 = 0} \\
\gl{2x - 5 = 11} & \gl{-x = 7} \\
\end{gleichungsraster}
\end{aufgabe}

\begin{aufgabe}{Ich kann durch Einsetzen prüfen, ob eine Zahl eine Lösung ist.}
Setze die Zahl ein und kreuze an: wahre Aussage (wA) oder falsche Aussage (fA).
\begin{teile}
\teil $x = 3$ \quad in \quad $x + 5 = 8$ \hfill \kreuz{wA} \kreuz{fA}
\teil $x = 2$ \quad in \quad $4x = 10$ \hfill \kreuz{wA} \kreuz{fA}
\teil $x = 0{,}5$ \quad in \quad $4x + 1 = 3$ \hfill \kreuz{wA} \kreuz{fA}
\teil $x = -2$ \quad in \quad $5 - x = 3$ \hfill \kreuz{wA} \kreuz{fA}
\teil $x = -3$ \quad in \quad $2x - 1 = -7$ \hfill \kreuz{wA} \kreuz{fA}
\end{teile}
\end{aufgabe}

\begin{aufgabe}{Ich kann Quadratwurzeln ziehen.}
Gib die Wurzel an. Gibt es sie nicht, schreibe „gibt es nicht“.
\begin{teilezwei}
\tz $\sqrt{49} = {}$ \leerfeld & \tz $\sqrt{100} = {}$ \leerfeld \\
\tz $\sqrt{0{,}64} = {}$ \leerfeld & \tz $\sqrt{196} = {}$ \leerfeld \\
\tz $\sqrt{-9} = {}$ \leerfeld & \\
\anweisung{Rechne mit dem Taschenrechner und runde auf zwei Stellen nach dem Komma.}
\tz $\sqrt{7} \approx {}$ \leerfeld & \\
\end{teilezwei}
\end{aufgabe}

\begin{aufgabe}{Ich kann den Scheitel aus der Scheitelpunktform ablesen.}
Lies den Scheitel S der Parabel ab.
\begin{teilezwei}
\tz $y = (x - 2)^2 + 1$ \quad \punktfeld & \tz $y = (x - 4)^2$ \quad \punktfeld \\
\tz $y = x^2 - 5$ \quad \punktfeld & \tz $y = (x + 3)^2 - 2$ \quad \punktfeld \\
\end{teilezwei}
\end{aufgabe}

\begin{aufgabe}{Ich kann Klammern ausmultiplizieren, auch mit den binomischen Formeln.}
Multipliziere aus.
\begin{gleichungsraster}[1]
\gl{2 \cdot (x + 5)} & \gl{x \cdot (x + 3)} \\
\gl{x \cdot (x - 6)} & \gl{(x + 6)^2} \\
\gl{(x - 9)^2} & \\
\end{gleichungsraster}
\end{aufgabe}

\begin{aufgabe}{Ich kann ausklammern.}
Klammere so viel wie möglich aus.
\begin{gleichungsraster}[1]
\gl{4x + 8} & \gl{x^2 + 3x} \\
\gl{x^2 - 9x} & \gl{2x^2 + 6x} \\
\gl{x^2 - x} & \\
\end{gleichungsraster}
\end{aufgabe}

\begin{aufgabe}{Ich kann Terme ordnen und zusammenfassen.}
Ordne nach $x^2$, $x$ und Zahl und fasse zusammen.
\begin{gleichungsraster}[1]
\gl{3 + x^2 + 2x} & \gl{x^2 + 4x + 3x + 1} \\
\gl{5x - 2 + x^2 - 6} & \gl{2x - x^2 + 4} \\
\gl{x^2 - 3x + 10 - 12 + x} & \\
\end{gleichungsraster}
\end{aufgabe}

\begin{aufgabe}{Ich finde den Fehler beim Auflösen einer quadrierten Klammer.}
Jonas rechnet so:
\rechnung{(x + 5)^2 - 30 &= x^2 + 25 - 30 \\ &= x^2 - 5}
Finde den Fehler, benenne ihn und korrigiere die Rechnung.
\begin{gleichungsraster}[3]
\gl{(x + 5)^2 - 30} & \\
\end{gleichungsraster}
\end{aufgabe}

\begin{aufgabe}{Ich kann eine quadrierte Klammer auflösen und zusammenfassen.}
Multipliziere aus und fasse zusammen.
\begin{gleichungsraster}[2]
\gl{(x + 2)^2 - 7} & \gl{(x - 3)^2 + 4x} \\
\gl{(x - 1)^2 - 2x + 5} & \\
\end{gleichungsraster}
\end{aufgabe}

\clearpage
% Einheit 1 – Wurzelziehen und Lösbarkeit – Hauptnummern 12 bis 23
\einheitenkopf[e1][1 Wurzelziehen]{Einheit 1 von 4 · Wurzelziehen und Lösbarkeit}
\zweigzeile{Hier lernst du, Gleichungen wie $x^2 = c$ durch Wurzelziehen zu lösen und zu sagen, ob es zwei, eine oder keine Lösung gibt · neu in diesem Jahr, je nach Buch bis Klasse 10 (am Gymnasium schon in Klasse 8 oder 9) · P10 · baut auf: Quadrieren, Wurzel ziehen, lineare Gleichungen, Scheitelpunktform}
\setcounter{aufgabe}{11}

\begin{aufgabe}{Ich erkenne, ob eine Gleichung quadratisch ist.}
Kreuze an: quadratisch oder linear? Löse nichts.
\begin{teile}
\teil $x^2 = 50$ \hfill \kreuz{quadratisch} \kreuz{linear}
\teil $5x - 3 = 12$ \hfill \kreuz{quadratisch} \kreuz{linear}
\teil $(x + 2)^2 = 1$ \hfill \kreuz{quadratisch} \kreuz{linear}
\teil $x \cdot x = 16$ \hfill \kreuz{quadratisch} \kreuz{linear}
\teil $3x + x = 12$ \hfill \kreuz{quadratisch} \kreuz{linear}
\end{teile}
\end{aufgabe}

\begin{aufgabe}{Ich erkenne, ob rechts eine positive Zahl, null oder eine negative Zahl steht.}
Sieh nur auf die rechte Seite und kreuze an. Rechne nichts.
\begin{teile}
\teil $x^2 = 64$ \hfill \kreuz{positiv} \kreuz{null} \kreuz{negativ}
\teil $x^2 = -4$ \hfill \kreuz{positiv} \kreuz{null} \kreuz{negativ}
\teil $2x^2 = 0$ \hfill \kreuz{positiv} \kreuz{null} \kreuz{negativ}
\teil $x^2 = 2{,}5$ \hfill \kreuz{positiv} \kreuz{null} \kreuz{negativ}
\teil $x^2 = -0{,}01$ \hfill \kreuz{positiv} \kreuz{null} \kreuz{negativ}
\end{teile}
\end{aufgabe}

\verfahren{Gleichungen der Form $x^2 = c$}
\begin{aufgabe}{Ich kann eine Gleichung $x^2 = c$ durch Wurzelziehen lösen.}
Löse die Gleichung. Gib beide Lösungen an.
\rechenplatz{2}
\begin{gleichungsraster}[2]
\gl{x^2 = 16} & \gl{x^2 = 81} \\
\gl{x^2 = 4} & \gl{x^2 = 144} \\
\anweisung{Lass die Wurzel stehen und gib zusätzlich Näherungswerte auf zwei Stellen nach dem Komma an.}
\gl{x^2 = 10} & \\
\anweisung{Hat die Gleichung keine Lösung, schreibe das als Satz.}
\gl{x^2 = -16} & \\
\end{gleichungsraster}
\end{aufgabe}

\begin{aufgabe}{Ich kann $x^2$ erst allein stellen und dann die Wurzel ziehen.}
Löse die Gleichung.
\rechenplatz{3}
\begin{gleichungsraster}[3]
\gl{x^2 - 4 = 12} & \gl{x^2 + 3 = 67} \\
\gl{2x^2 = 50} & \gl{3x^2 + 1 = 13} \\
\gl{2x^2 + 9 = 1} & \gl{4x^2 + 3 = 3} \\
\end{gleichungsraster}
\end{aufgabe}

\verfahren{Gleichungen mit quadrierter Klammer}
\begin{aufgabe}{Ich kann eine Gleichung mit quadrierter Klammer rückwärts lösen.}
Löse die Gleichung.
\rechenplatz{3}
\begin{gleichungsraster}[3]
\gl{(x - 2)^2 = 9} & \gl{(x - 1)^2 = 25} \\
\gl{(x - 5)^2 = 4} & \gl{(x - 3)^2 = 1} \\
\gl{(x + 2)^2 = 25} & \gl{(x + 5)^2 = 0} \\
\end{gleichungsraster}
\end{aufgabe}

\verfahren{Lösungen prüfen und zählen}
\begin{aufgabe}{Ich kann prüfen, ob eine Zahl eine quadratische Gleichung löst.}
Setze die Zahl ein und kreuze an: wahre Aussage (wA) oder falsche Aussage (fA).
\begin{teile}
\teil $x = 7$ \quad in \quad $x^2 + 1 = 50$ \hfill \kreuz{wA} \kreuz{fA}
\teil $x = -7$ \quad in \quad $x^2 + 1 = 50$ \hfill \kreuz{wA} \kreuz{fA}
\teil $x = 3$ \quad in \quad $(x - 1)^2 = 9$ \hfill \kreuz{wA} \kreuz{fA}
\teil $x = -2$ \quad in \quad $(x + 5)^2 = 9$ \hfill \kreuz{wA} \kreuz{fA}
\teil $x = -1$ \quad in \quad $3x^2 - 5 = -8$ \hfill \kreuz{wA} \kreuz{fA}
\end{teile}
\end{aufgabe}

\begin{aufgabe}{Ich kann am Graphen ablesen, wie viele Lösungen eine Gleichung hat.}
Die Parabel gehört zu $y = (x + 1)^2$. Lies für jede Gleichung ab, wie viele Lösungen sie hat, und gib die Lösungen an.

\begin{ksys}[xmin=-5,xmax=4,ymin=-3,ymax=6,ablesen]
\parabel{1}{-1}{0}{}
\gerade[3]{0}{4}{y=4}
\gerade[3]{0}{1}{y=1}
\gerade[3]{0}{-2}{y=-2}
\end{ksys}

\begin{teile}
\teil $(x + 1)^2 = 4$ \hfill Anzahl: \leerfeld \quad Lösungen: \leerfeld
\teil $(x + 1)^2 = 1$ \hfill Anzahl: \leerfeld \quad Lösungen: \leerfeld
\teil $(x + 1)^2 = 0$ \hfill Anzahl: \leerfeld \quad Lösungen: \leerfeld
\teil $(x + 1)^2 = -2$ \hfill Anzahl: \leerfeld \quad Lösungen: \leerfeld
\end{teile}
\end{aufgabe}

\begin{aufgabe}{Ich kann Gleichungen mit keiner, einer oder zwei Lösungen angeben.}
Ergänze die rechte Seite so, dass die Gleichung die angegebene Zahl von Lösungen hat.
\begin{teile}
\teil genau eine Lösung: \quad $(x - 3)^2 = {}$ \leerfeld
\teil zwei Lösungen: \quad $(x + 2)^2 = {}$ \leerfeld
\teil keine Lösung: \quad $x^2 + 5 = {}$ \leerfeld
\teil keine Lösung: \quad $3x^2 = {}$ \leerfeld
\anweisung{Schreibe selbst eine Gleichung.}
\teil eine Gleichung mit quadrierter Klammer, die keine Lösung hat: \leerfeld
\end{teile}
\end{aufgabe}

\begin{aufgabe}{Ich kann Aussagen über die Lösungen einer Gleichung beurteilen.}
Kreuze an, ob die Aussage wahr oder falsch ist.
\begin{teile}
\teil \sachtabelle{lcc}{Aussage & wahr & falsch}{„$(x - 6)^2 = 0$ hat genau eine Lösung.“ & $\square$ & $\square$ \\ „$2$ und $10$ sind Lösungen von $(x + 6)^2 = 16$.“ & $\square$ & $\square$}
\teil Gib eine quadratische Gleichung an, die keine Lösung hat: \leerfeld \hfill (P10 2021 OS)
\end{teile}
\end{aufgabe}

\begin{aufgabe}{Ich finde den Fehler beim Wurzelziehen.}
Finde den Fehler, benenne ihn und korrigiere die Rechnung.
\begin{teile}
\teil Paul rechnet: \rechnung{x^2 &= 64 \\ x &= 8}
\teil Lea rechnet: \rechnung{(x + 9)^2 &= 4 \\ x + 9 &= 2 \quad \text{oder} \quad x + 9 = -2 \\ x_1 &= 11 \quad \text{oder} \quad x_2 = 7}
\teil Ben rechnet: \rechnung{x^2 &= -81 \\ x &= -9}
\end{teile}
\end{aufgabe}

\begin{aufgabe}{Ich kann begründen, wie viele Lösungen eine Gleichung hat.}
Begründe ohne zu rechnen.
\begin{teile}
\teil Warum hat $x^2 = 20$ zwei Lösungen?
\teil Warum hat $x^2 = -20$ keine Lösung?
\teil Warum hat $(x - 9)^2 = 0$ nur eine Lösung?
\end{teile}
\end{aufgabe}

\begin{aufgabe}{Ich kann mit Wurzelziehen eine Länge berechnen.}
Berechne. Runde auf zwei Stellen nach dem Komma.
\begin{teile}
\teil Ein quadratisches Beet hat eine Fläche von $121\,\text{m}^2$. Wie lang ist eine Seite? \\ Seitenlänge: \leerfeld[m]
\teil Eine quadratische Tischplatte hat eine Fläche von $1{,}44\,\text{m}^2$. Wie lang ist eine Seite? \\ Seitenlänge: \leerfeld[m]
\teil Eine zylinderförmige Dose soll $10\,\text{cm}$ hoch sein und $500\,\text{cm}^3$ fassen. Für das Volumen gilt $V = \pi \cdot r^2 \cdot h$. Wie groß muss der Radius $r$ sein? (P10 2022 OS) \\ $r = {}$ \leerfeld[cm]
\end{teile}

\zylinder{1.2}{2}{r}{h}

\end{aufgabe}

\clearpage
% Einheit 2 – Normalform und p-q-Formel – Hauptnummern 24 bis 39
\einheitenkopf[e2][2 p-q-Formel]{Einheit 2 von 4 · Normalform und p-q-Formel}
\zweigzeile{Hier lernst du, jede quadratische Gleichung zu ordnen und mit der p-q-Formel zu lösen · neu in diesem Jahr, je nach Buch bis Klasse 10 · P10 oft · baut auf: Wurzelziehen (Einheit 1), Ausmultiplizieren, Terme ordnen}
\setcounter{aufgabe}{23}

\begin{aufgabe}{Ich erkenne an der Form einer Gleichung, welcher Lösungsweg passt.}
Kreuze den passenden Lösungsweg an. Löse nichts.
\begin{teile}
\teil $x^2 - 7x + 10 = 0$ \hfill \kreuz{Wurzelziehen} \kreuz{rückwärts rechnen} \kreuz{p-q-Formel}
\teil $x^2 = 30$ \hfill \kreuz{Wurzelziehen} \kreuz{rückwärts rechnen} \kreuz{p-q-Formel}
\teil $(x - 4)^2 = 5$ \hfill \kreuz{Wurzelziehen} \kreuz{rückwärts rechnen} \kreuz{p-q-Formel}
\teil $3x^2 - 12 = 0$ \hfill \kreuz{Wurzelziehen} \kreuz{rückwärts rechnen} \kreuz{p-q-Formel}
\teil $x^2 + x = 6$ \hfill \kreuz{Wurzelziehen} \kreuz{rückwärts rechnen} \kreuz{p-q-Formel}
\end{teile}
\end{aufgabe}

\begin{aufgabe}{Ich erkenne p und q in einer Gleichung in Normalform.}
Kreise p ein und unterstreiche q, jeweils mit dem Vorzeichen davor. Rechne nichts.
\begin{teilezwei}
\tz $x^2 + 7x + 2 = 0$ & \tz $x^2 - 2x + 9 = 0$ \\
\tz $x^2 + 5x - 1 = 0$ & \tz $x^2 - x - 20 = 0$ \\
\tz $x^2 - 11 = 0$ & \\
\end{teilezwei}
\end{aufgabe}

\begin{aufgabe}{Ich erkenne, ob vor $x^2$ eine Eins steht.}
Kreuze an, ob du vor der p-q-Formel teilen musst. Unterstreiche die Zahl, durch die du teilen musst.
\begin{teile}
\teil $x^2 + 5x - 3 = 0$ \hfill \kreuz{teilen} \kreuz{nicht teilen}
\teil $2x^2 + 4x - 7 = 0$ \hfill \kreuz{teilen} \kreuz{nicht teilen}
\teil $-x^2 + 3x + 4 = 0$ \hfill \kreuz{teilen} \kreuz{nicht teilen}
\teil $0{,}5x^2 - x + 2 = 0$ \hfill \kreuz{teilen} \kreuz{nicht teilen}
\teil $x^2 - 9x + 1 = 0$ \hfill \kreuz{teilen} \kreuz{nicht teilen}
\end{teile}
\end{aufgabe}

\verfahren{Lösen mit der p-q-Formel}
\begin{aufgabe}{Ich kann eine Gleichung in Normalform mit der p-q-Formel lösen. (je nach Buch erst in Klasse 10)}
Löse mit der p-q-Formel.
\rechenplatz{4}
\begin{gleichungsraster}[3]
\gl{x^2 + 6x + 8 = 0} & \gl{x^2 + 10x + 16 = 0} \\
\gl{x^2 + 6x + 5 = 0} & \gl{x^2 + 10x + 21 = 0} \\
\end{gleichungsraster}
\end{aufgabe}

\begin{aufgabe}{Ich kann eine Gleichung in Normalform mit der p-q-Formel lösen – weiter.}
\begin{gleichungsraster}[3]
\gl{x^2 - 4x - 12 = 0} & \gl{x^2 + 3x - 10 = 0} \\
\anweisung{Gib die Lösungen mit Wurzel und als Näherungswerte auf zwei Stellen nach dem Komma an.}
\gl{x^2 - 4x + 1 = 0} & \gl{x^2 + x - 3 = 0} \\
\anweisung{Berechne die Nullstellen der Parabel. (P10 2025 OS)}
\gl{y = x^2 - 8x + 14} & \\
\end{gleichungsraster}
\end{aufgabe}

\verfahren{Gleichungen ordnen und durch die Zahl vor $x^2$ teilen}
\begin{aufgabe}{Ich kann eine Gleichung erst ordnen und dann mit der p-q-Formel lösen.}
Löse die Gleichung.
\rechenplatz{5}
\begin{gleichungsraster}[3]
\gl{x^2 + 2x = 8} & \gl{x^2 = 6x - 5} \\
\gl{x^2 + 4 = 5x} & \gl{x^2 + 9x - 5 = 3x + 11} \\
\end{gleichungsraster}
\end{aufgabe}

\begin{aufgabe}{Ich kann eine Gleichung mit einer Zahl vor $x^2$ lösen.}
Löse die Gleichung.
\rechenplatz{5}
\begin{gleichungsraster}[3]
\gl{2x^2 - 4x - 6 = 0} & \gl{3x^2 + 9x - 12 = 0} \\
\gl{-x^2 + 6x - 8 = 0} & \gl{0{,}5x^2 - x - 4 = 0} \\
\anweisung{Berechne die Nullstellen der Parabel. (P10 2020 OS)}
\gl{y = 3x^2 - 3x - 36} & \\
\end{gleichungsraster}
\end{aufgabe}

\begin{aufgabe}{Ich kann eine Gleichung mit Klammern erst auflösen, ordnen und dann lösen.}
Löse die Gleichung.
\rechenplatz{5}
\begin{gleichungsraster}[3]
\gl{x \cdot (x + 4) = 12} & \gl{(x + 1)^2 = 3x + 7} \\
\gl{(x - 3)^2 - 1 = -2x + 8} & \\
\anweisung{Für welche x hat die Parabel den y-Wert $-8$? (P10 2023 OS)}
\gl{y = -x^2 + 6x - 1} & \\
\end{gleichungsraster}
\end{aufgabe}

\verfahren{Zahl der Lösungen}
\begin{aufgabe}{Ich kann eine Gleichung mit genau einer Lösung lösen.}
Löse mit der p-q-Formel.
\begin{gleichungsraster}[3]
\gl{x^2 + 6x + 9 = 0} & \gl{x^2 - 10x + 25 = 0} \\
\gl{x^2 - x + 0{,}25 = 0} & \gl{4x^2 + 12x + 9 = 0} \\
\end{gleichungsraster}
\end{aufgabe}

\begin{aufgabe}{Ich kann erkennen, dass eine Gleichung keine Lösung hat.}
Löse mit der p-q-Formel. Hat die Gleichung keine Lösung, schreibe das als Satz.
\begin{gleichungsraster}[3]
\gl{x^2 + 2x + 5 = 0} & \gl{x^2 - 6x + 10 = 0} \\
\gl{x^2 + x + 1 = 0} & \gl{2x^2 + 4x + 7 = 0} \\
\end{gleichungsraster}
\end{aufgabe}

\begin{aufgabe}{Ich kann am Wert unter der Wurzel entscheiden, wie viele Lösungen es gibt.}
Berechne nur den Wert unter der Wurzel und kreuze an.
\begin{teile}
\teil $x^2 + 4x + 1 = 0$ \hfill Wert: \leerfeld \quad \kreuz{zwei} \kreuz{eine} \kreuz{keine}
\teil $x^2 - 12x + 36 = 0$ \hfill Wert: \leerfeld \quad \kreuz{zwei} \kreuz{eine} \kreuz{keine}
\teil $x^2 + 2x + 4 = 0$ \hfill Wert: \leerfeld \quad \kreuz{zwei} \kreuz{eine} \kreuz{keine}
\teil $x^2 - 3x + 1 = 0$ \hfill Wert: \leerfeld \quad \kreuz{zwei} \kreuz{eine} \kreuz{keine}
\teil $x^2 + 5x + 7 = 0$ \hfill Wert: \leerfeld \quad \kreuz{zwei} \kreuz{eine} \kreuz{keine}
\anweisung{Den Wert unter der Wurzel nennt man auch Diskriminante.}
\teil $x^2 - 2x - 1 = 0$ \hfill Diskriminante: \leerfeld \quad \kreuz{zwei} \kreuz{eine} \kreuz{keine}
\end{teile}
\end{aufgabe}

\verfahren{Probe und Lösungsweg}
\begin{aufgabe}{Ich kann mit einer Probe prüfen, ob eine Lösung stimmt.}
Setze die Zahl ein und kreuze an: wahre Aussage (wA) oder falsche Aussage (fA).
\begin{teile}
\teil $x = 2$ \quad in \quad $x^2 + 2x - 8 = 0$ \hfill \kreuz{wA} \kreuz{fA}
\teil $x = -1$ \quad in \quad $x^2 - 5x + 4 = 0$ \hfill \kreuz{wA} \kreuz{fA}
\teil $x = -1$ \quad in \quad $x^2 - 2x - 3 = 0$ \hfill \kreuz{wA} \kreuz{fA}
\teil $x = 3$ \quad in \quad $x^2 + x - 6 = 0$ \hfill \kreuz{wA} \kreuz{fA}
\teil $x = -4$ \quad in \quad $x^2 + 3x - 4 = 0$ \hfill \kreuz{wA} \kreuz{fA}
\end{teile}
\end{aufgabe}

\begin{aufgabe}{Ich kann den kürzesten Lösungsweg wählen.}
Löse die Gleichung auf dem kürzesten Weg.
\begin{gleichungsraster}[3]
\gl{x^2 - 20 = 5} & \gl{x^2 + 2x - 24 = 0} \\
\gl{(x + 3)^2 = 16} & \gl{x^2 - 6x = 16} \\
\end{gleichungsraster}
\end{aufgabe}

\begin{aufgabe}{Ich finde den Fehler beim Lösen mit der p-q-Formel.}
Finde den Fehler, benenne ihn und korrigiere die Rechnung.
\begin{teile}
\teil Mia rechnet: \rechnung{2x^2 + 16x + 14 &= 0 \qquad p = 16,\ q = 14 \\ x_{1,2} &= -8 \pm \sqrt{64 - 14} = -8 \pm \sqrt{50}}
\teil Tim rechnet: \rechnung{x^2 - 10x + 9 &= 0 \\ x_{1,2} &= -5 \pm \sqrt{25 - 9} = -5 \pm 4 \\ x_1 &= -1, \quad x_2 = -9}
\teil Ella rechnet: \rechnung{x^2 - 8x + 12 &= 0 \\ x_{1,2} &= 4 \pm \sqrt{16 + 12} = 4 \pm \sqrt{28}}
\end{teile}
\end{aufgabe}

\begin{aufgabe}{Ich kann begründen, wann die p-q-Formel funktioniert.}
Begründe.
\begin{teile}
\teil Warum darf man die p-q-Formel bei $2x^2 + 6x - 8 = 0$ nicht sofort anwenden?
\teil Warum hat eine Gleichung keine Lösung, wenn der Wert unter der Wurzel negativ ist?
\end{teile}
\end{aufgabe}

\begin{aufgabe}{Ich kann berechnen, wo sich eine Gerade und eine Parabel schneiden.}
Gezeichnet sind die Parabel zu $y = x^2 - 2x - 1$ und die Gerade zu $y = x + 3$.

\begin{ksys}[xmin=-3,xmax=6,ymin=-3,ymax=8,ablesen]
\parabel{1}{1}{-2}{}
\gerade[5]{1}{3}{}
\end{ksys}

\begin{teile}
\teil Setze gleich und berechne die x-Werte der Schnittpunkte.
\teil Berechne die y-Werte und gib die Schnittpunkte an.
\teil Kontrolliere deine Schnittpunkte am Graphen. \hfill (P10 2024 OS)
\end{teile}
\end{aufgabe}

\clearpage
% Einheit 3 – Sachaufgaben – Hauptnummern 40 bis 46
\einheitenkopf[e3][3 Sachaufgaben]{Einheit 3 von 4 · Sachaufgaben}
\zweigzeile{Hier lernst du, Zahlenrätsel und Flächenaufgaben mit quadratischen Gleichungen zu lösen · neu in diesem Jahr, je nach Buch bis Klasse 10 · keine P10-Aufgabe · baut auf: Wurzelziehen (Einheit 1), p-q-Formel (Einheit 2)}
\setcounter{aufgabe}{39}

\begin{aufgabe}{Ich erkenne, welche Lösung zur Frage passt.}
Die Gleichung ist schon gelöst. Kreuze an, welche Lösung als Antwort passt. Rechne nichts.
\begin{teile}
\teil Gesucht ist die Seitenlänge eines Quadrats. Lösungen: $9$ und $-9$ \hfill \kreuz{$9$} \kreuz{$-9$} \kreuz{beide}
\teil Gesucht ist eine Zahl. Lösungen: $4$ und $-4$ \hfill \kreuz{$4$} \kreuz{$-4$} \kreuz{beide}
\teil Gesucht ist die Breite eines Rechtecks. Lösungen: $6$ und $-10$ \hfill \kreuz{$6$} \kreuz{$-10$} \kreuz{beide}
\teil Gesucht ist die Zahl der Sitzreihen im Kino. Lösungen: $-15$ und $12$ \hfill \kreuz{$-15$} \kreuz{$12$} \kreuz{beide}
\teil Gesucht ist eine natürliche Zahl. Lösungen: $7$ und $-8$ \hfill \kreuz{$7$} \kreuz{$-8$} \kreuz{beide}
\end{teile}
\end{aufgabe}

\verfahren{Zahlenrätsel}
\begin{aufgabe}{Ich kann ein Zahlenrätsel mit dem Quadrat einer Zahl lösen.}
Löse die Gleichung und gib alle passenden Zahlen an.
\begin{teile}
\teil Das Quadrat einer Zahl ist 169. \quad $x^2 = 169$ \hfill Zahlen: \leerfeld
\teil Das Quadrat einer Zahl ist 400. \quad $x^2 = 400$ \hfill Zahlen: \leerfeld
\teil Verdoppelt man das Quadrat einer Zahl, erhält man 128. \quad $2x^2 = 128$ \hfill Zahlen: \leerfeld
\teil Das Quadrat einer Zahl, vermindert um 10, ist 90. \quad $x^2 - 10 = 90$ \hfill Zahlen: \leerfeld
\end{teile}
\end{aufgabe}

\begin{aufgabe}{Ich kann zu einem Zahlenrätsel selbst die Gleichung aufstellen.}
Stelle eine Gleichung auf, löse sie und gib die passenden Zahlen an.
\begin{teile}
\teil Addiert man zum Quadrat einer Zahl 15, erhält man 96. \\ Gleichung: \leerfeld \quad Zahlen: \leerfeld
\teil Das Dreifache des Quadrats einer Zahl ist 75. \\ Gleichung: \leerfeld \quad Zahlen: \leerfeld
\teil Das Produkt einer Zahl mit ihrem Nachfolger ist 56. \\ Gleichung: \leerfeld \quad Zahlen: \leerfeld
\teil Das Produkt zweier aufeinanderfolgender natürlicher Zahlen ist 132. \\ Gleichung: \leerfeld \quad Zahlen: \leerfeld
\end{teile}
\end{aufgabe}

\verfahren{Flächenaufgaben}
\begin{aufgabe}{Ich kann eine Rechteckaufgabe mit einer quadratischen Gleichung lösen.}
Ein rechteckiger Spielplatz ist $3\,\text{m}$ länger als breit. Seine Fläche beträgt $40\,\text{m}^2$.

\rechteck[seiten={x+3,x,x+3,x},punkte=]{4}{2.5}

\begin{teile}
\teil Stelle eine Gleichung für die Fläche auf.
\teil Löse die Gleichung.
\teil Welche Lösung passt zur Frage? Begründe.
\teil Gib Breite und Länge des Spielplatzes an und mache die Probe.
\end{teile}
\end{aufgabe}

\begin{aufgabe}{Ich kann eine Flächenaufgabe mit Klammern lösen.}
Stelle eine Gleichung auf, löse sie und antworte mit Einheit.
\begin{teile}
\teil Ein quadratisches Foto steckt in einem Rahmen, der ringsum $2\,\text{cm}$ breit ist. Foto und Rahmen zusammen haben eine Fläche von $196\,\text{cm}^2$. Wie lang ist eine Seite des Fotos?
\end{teile}

\quadratrand{x}{2 cm}

\begin{teile}
\teil Ein quadratischer Garten wird auf einer Seite um $4\,\text{m}$ verlängert und auf der anderen Seite um $2\,\text{m}$ verkürzt. Das neue Rechteck hat eine Fläche von $55\,\text{m}^2$. Wie lang war eine Seite des Gartens?
\end{teile}

\rechteck[seiten={x+4,x-2,x+4,x-2},punkte=]{4}{2}

\end{aufgabe}

\begin{aufgabe}{Ich finde den Fehler in einer Sachaufgabe.}
Finde den Fehler, benenne ihn und korrigiere die Rechnung.
\begin{teile}
\teil Ein Rechteck ist $5\,\text{cm}$ länger als breit und hat eine Fläche von $24\,\text{cm}^2$. Tim rechnet: \rechnung{x \cdot (x + 5) &= 24 \\ x^2 + 5x - 24 &= 0 \\ x_1 &= 3, \quad x_2 = -8} Tims Antwort: „Die Breite ist $3\,\text{cm}$ oder $-8\,\text{cm}$.“
\teil Ein Rechteck ist $2\,\text{cm}$ länger als breit und hat eine Fläche von $48\,\text{cm}^2$. Mia rechnet: \rechnung{2x + 2 \cdot (x + 2) &= 48 \\ 4x + 4 &= 48 \\ x &= 11}
\end{teile}
\end{aufgabe}

\begin{aufgabe}{Ich kann begründen, welche Lösungen im Sachzusammenhang gelten.}
Begründe.
\begin{teile}
\teil Warum gilt bei der Seitenlänge eines Rechtecks nur die positive Lösung?
\teil Warum sind beim Rätsel „Das Quadrat einer Zahl ist 225.“ beide Lösungen richtig?
\end{teile}
\end{aufgabe}

\clearpage
% Ausblick Einheit 4 – Satz vom Nullprodukt – Hauptnummern 47 bis 55
\einheitenkopf[e4][Ausblick 4 Nullprodukt]{Ausblick Einheit 4 von 4 · Satz vom Nullprodukt}
\zweigzeile{Hier lernst du, Gleichungen in Produktform mit dem Satz vom Nullprodukt zu lösen · kommt nächstes Jahr (am Gymnasium schon in Klasse 8 oder 9) · P10 · baut auf: Ausklammern, lineare Gleichungen, p-q-Formel (Einheit 2)}
\setcounter{aufgabe}{46}

\begin{aufgabe}{Ich erkenne, ob rechts eine Null steht.}
Kreuze an, ob der Satz vom Nullprodukt gilt, weil rechts null steht. Löse nichts.
\begin{teile}
\teil $(x - 3) \cdot (x + 1) = 0$ \hfill \kreuz{gilt} \kreuz{gilt nicht}
\teil $x \cdot (x + 2) = 8$ \hfill \kreuz{gilt} \kreuz{gilt nicht}
\teil $(x + 4) \cdot (x - 4) = 0$ \hfill \kreuz{gilt} \kreuz{gilt nicht}
\teil $x \cdot (x - 5) = -4$ \hfill \kreuz{gilt} \kreuz{gilt nicht}
\teil $2x \cdot (x - 1) = 0$ \hfill \kreuz{gilt} \kreuz{gilt nicht}
\end{teile}
\end{aufgabe}

\verfahren{Produkt gleich null}
\begin{aufgabe}{Ich kann eine Gleichung in Produktform mit dem Satz vom Nullprodukt lösen.}
Löse die Gleichung.
\rechenplatz{2}
\begin{gleichungsraster}[2]
\gl{(x - 2) \cdot (x - 5) = 0} & \gl{(x - 1) \cdot (x - 6) = 0} \\
\gl{(x - 3) \cdot (x - 4) = 0} & \gl{(x - 7) \cdot (x - 2) = 0} \\
\gl{(x + 5) \cdot (x - 1) = 0} & \gl{x \cdot (x + 6) = 0} \\
\gl{(x - 8) \cdot (x - 8) = 0} & \\
\end{gleichungsraster}
\end{aufgabe}

\begin{aufgabe}{Ich kann prüfen, welche Zahl eine Gleichung in Produktform erfüllt.}
Setze die Zahl ein und kreuze an: wahre Aussage (wA) oder falsche Aussage (fA).
\begin{teile}
\teil $x = 3$ \quad in \quad $(x - 3) \cdot (x + 2) = 0$ \hfill \kreuz{wA} \kreuz{fA}
\teil $x = -2$ \quad in \quad $x \cdot (x + 4) = -4$ \hfill \kreuz{wA} \kreuz{fA}
\teil $x = 2$ \quad in \quad $x \cdot (x + 4) = -4$ \hfill \kreuz{wA} \kreuz{fA}
\teil $x = -1$ \quad in \quad $(x + 3) \cdot (x - 2) = -6$ \hfill \kreuz{wA} \kreuz{fA}
\teil $x = 1$ \quad in \quad $(x - 4) \cdot (x + 2) = 9$ \hfill \kreuz{wA} \kreuz{fA}
\anweisung{Kreuze den Wert an, der die Gleichung erfüllt. (P10 2025 OS)}
\teil $x \cdot (x + 7) = -10$ \hfill \kreuz{$x = 2$} \kreuz{$x = 5$} \kreuz{$x = -2$} \kreuz{$x = -10$}
\end{teile}
\end{aufgabe}

\verfahren{Ausklammern}
\begin{aufgabe}{Ich kann eine Gleichung ohne Zahlglied durch Ausklammern lösen.}
Löse die Gleichung.
\rechenplatz{3}
\begin{gleichungsraster}[3]
\gl{x^2 + 4x = 0} & \gl{x^2 + 9x = 0} \\
\gl{x^2 - x = 0} & \gl{3x^2 - 12x = 0} \\
\gl{x^2 = 5x} & \\
\end{gleichungsraster}
\end{aufgabe}

\verfahren{Produkt ungleich null}
\begin{aufgabe}{Ich kann eine Produktgleichung lösen, bei der rechts nicht null steht.}
Löse die Gleichung.
\rechenplatz{4}
\begin{gleichungsraster}[3]
\gl{x \cdot (x + 2) = 35} & \gl{(x - 2) \cdot (x + 3) = 14} \\
\gl{x \cdot (x - 3) = 10} & \\
\end{gleichungsraster}
\end{aufgabe}

\begin{aufgabe}{Ich kann eine Gleichung lösen, indem ich die linke Seite als Produkt schreibe.}
Schreibe die linke Seite als Produkt zweier Klammern und löse.
\begin{gleichungsraster}[2]
\gl{x^2 + 8x + 15 = 0} & \gl{x^2 - 7x + 12 = 0} \\
\gl{x^2 + 2x - 3 = 0} & \\
\anweisung{Berechne die Nullstellen der Parabel durch Faktorisieren. (P10 2020 OS)}
\gl{y = 2x^2 + 14x + 20} & \\
\end{gleichungsraster}
\end{aufgabe}

\begin{aufgabe}{Ich finde den Fehler beim Satz vom Nullprodukt.}
Finde den Fehler, benenne ihn und korrigiere die Rechnung.
\begin{teile}
\teil Nora rechnet: \rechnung{x^2 &= 3x &&\mid : x \\ x &= 3}
\teil Ali rechnet: \rechnung{(x + 2) \cdot (x - 6) &= 0 \\ x_1 &= 2, \quad x_2 = -6}
\teil Kim rechnet: \rechnung{x \cdot (x + 1) &= 12 \\ x = 12 \quad &\text{oder} \quad x + 1 = 12 \\ x_1 = 12 \quad &\text{oder} \quad x_2 = 11}
\end{teile}
\end{aufgabe}

\begin{aufgabe}{Ich kann begründen, wann der Satz vom Nullprodukt gilt.}
Begründe.
\begin{teile}
\teil Warum ist $(x - 4) \cdot (x + 1)$ gleich null, wenn $x = 4$ ist?
\teil Warum darf man $x^2 = 6x$ nicht durch $x$ teilen?
\end{teile}
\end{aufgabe}

\begin{aufgabe}{Ich kann mit dem Satz vom Nullprodukt eine Sachfrage beantworten.}
Ein Ball wird vom Boden senkrecht nach oben geworfen. Nach $t$ Sekunden ist er $h = 15t - 5t^2$ Meter hoch.
\begin{teile}
\teil Schreibe $15t - 5t^2$ als Produkt.
\teil Nach wie vielen Sekunden ist der Ball wieder am Boden?
\teil Was bedeutet die Lösung $t = 0$ im Sachzusammenhang?
\end{teile}
\end{aufgabe}

% Abhakseite „Das kann ich" – Zone nur im Gesamt (\mitzone)
\begin{abhakseite}
\ifdefined\mitzone
\abhakgruppe{Kennst du schon}
\abhak{1}{Ich kann Zahlen quadrieren, auch negative Zahlen und Dezimalzahlen.}
\abhak{2}{Ich kann mit Klammern und negativen Zahlen rechnen: Klammer zuerst, Punkt vor Strich.}
\abhak{3}{Ich kann eine lineare Gleichung lösen.}
\abhak{4}{Ich kann durch Einsetzen prüfen, ob eine Zahl eine Lösung ist.}
\abhak{5}{Ich kann Quadratwurzeln ziehen.}
\abhak{6}{Ich kann den Scheitel aus der Scheitelpunktform ablesen.}
\abhak{7}{Ich kann Klammern ausmultiplizieren, auch mit den binomischen Formeln.}
\abhak{8}{Ich kann ausklammern.}
\abhak{9}{Ich kann Terme ordnen und zusammenfassen.}
\abhak{10}{Ich finde den Fehler beim Auflösen einer quadrierten Klammer.}
\abhak{11}{Ich kann eine quadrierte Klammer auflösen und zusammenfassen.}
\fi
\abhakgruppe{Einheit 1 · Wurzelziehen und Lösbarkeit}
\abhak{12}{Ich erkenne, ob eine Gleichung quadratisch ist.}
\abhak{13}{Ich erkenne, ob rechts eine positive Zahl, null oder eine negative Zahl steht.}
\abhakgruppe{Einheit 1 · Gleichungen der Form $x^2 = c$}
\abhak{14}{Ich kann eine Gleichung $x^2 = c$ durch Wurzelziehen lösen.}
\abhak{15}{Ich kann $x^2$ erst allein stellen und dann die Wurzel ziehen.}
\abhakgruppe{Einheit 1 · Gleichungen mit quadrierter Klammer}
\abhak{16}{Ich kann eine Gleichung mit quadrierter Klammer rückwärts lösen.}
\abhakgruppe{Einheit 1 · Lösungen prüfen und zählen}
\abhak{17}{Ich kann prüfen, ob eine Zahl eine quadratische Gleichung löst.}
\abhak{18}{Ich kann am Graphen ablesen, wie viele Lösungen eine Gleichung hat.}
\abhak{19}{Ich kann Gleichungen mit keiner, einer oder zwei Lösungen angeben.}
\abhak{20}{Ich kann Aussagen über die Lösungen einer Gleichung beurteilen.}
\abhak{21}{Ich finde den Fehler beim Wurzelziehen.}
\abhak{22}{Ich kann begründen, wie viele Lösungen eine Gleichung hat.}
\abhak{23}{Ich kann mit Wurzelziehen eine Länge berechnen.}
\abhakgruppe{Einheit 2 · Normalform und p-q-Formel}
\abhak{24}{Ich erkenne an der Form einer Gleichung, welcher Lösungsweg passt.}
\abhak{25}{Ich erkenne p und q in einer Gleichung in Normalform.}
\abhak{26}{Ich erkenne, ob vor $x^2$ eine Eins steht.}
\abhakgruppe{Einheit 2 · Lösen mit der p-q-Formel}
\abhak{27}{Ich kann eine Gleichung in Normalform mit der p-q-Formel lösen. (je nach Buch erst in Klasse 10)}
\abhak{28}{Ich kann eine Gleichung in Normalform mit der p-q-Formel lösen – weiter.}
\abhakgruppe{Einheit 2 · Gleichungen ordnen und durch die Zahl vor $x^2$ teilen}
\abhak{29}{Ich kann eine Gleichung erst ordnen und dann mit der p-q-Formel lösen.}
\abhak{30}{Ich kann eine Gleichung mit einer Zahl vor $x^2$ lösen.}
\abhak{31}{Ich kann eine Gleichung mit Klammern erst auflösen, ordnen und dann lösen.}
\abhakgruppe{Einheit 2 · Zahl der Lösungen}
\abhak{32}{Ich kann eine Gleichung mit genau einer Lösung lösen.}
\abhak{33}{Ich kann erkennen, dass eine Gleichung keine Lösung hat.}
\abhak{34}{Ich kann am Wert unter der Wurzel entscheiden, wie viele Lösungen es gibt.}
\abhakgruppe{Einheit 2 · Probe und Lösungsweg}
\abhak{35}{Ich kann mit einer Probe prüfen, ob eine Lösung stimmt.}
\abhak{36}{Ich kann den kürzesten Lösungsweg wählen.}
\abhak{37}{Ich finde den Fehler beim Lösen mit der p-q-Formel.}
\abhak{38}{Ich kann begründen, wann die p-q-Formel funktioniert.}
\abhak{39}{Ich kann berechnen, wo sich eine Gerade und eine Parabel schneiden.}
\abhakgruppe{Einheit 3 · Sachaufgaben}
\abhak{40}{Ich erkenne, welche Lösung zur Frage passt.}
\abhakgruppe{Einheit 3 · Zahlenrätsel}
\abhak{41}{Ich kann ein Zahlenrätsel mit dem Quadrat einer Zahl lösen.}
\abhak{42}{Ich kann zu einem Zahlenrätsel selbst die Gleichung aufstellen.}
\abhakgruppe{Einheit 3 · Flächenaufgaben}
\abhak{43}{Ich kann eine Rechteckaufgabe mit einer quadratischen Gleichung lösen.}
\abhak{44}{Ich kann eine Flächenaufgabe mit Klammern lösen.}
\abhak{45}{Ich finde den Fehler in einer Sachaufgabe.}
\abhak{46}{Ich kann begründen, welche Lösungen im Sachzusammenhang gelten.}
\abhakgruppe{Ausblick Einheit 4 · Satz vom Nullprodukt}
\abhak{47}{Ich erkenne, ob rechts eine Null steht.}
\abhakgruppe{Ausblick Einheit 4 · Produkt gleich null}
\abhak{48}{Ich kann eine Gleichung in Produktform mit dem Satz vom Nullprodukt lösen.}
\abhak{49}{Ich kann prüfen, welche Zahl eine Gleichung in Produktform erfüllt.}
\abhakgruppe{Ausblick Einheit 4 · Ausklammern}
\abhak{50}{Ich kann eine Gleichung ohne Zahlglied durch Ausklammern lösen.}
\abhakgruppe{Ausblick Einheit 4 · Produkt ungleich null}
\abhak{51}{Ich kann eine Produktgleichung lösen, bei der rechts nicht null steht.}
\abhak{52}{Ich kann eine Gleichung lösen, indem ich die linke Seite als Produkt schreibe.}
\abhak{53}{Ich finde den Fehler beim Satz vom Nullprodukt.}
\abhak{54}{Ich kann begründen, wann der Satz vom Nullprodukt gilt.}
\abhak{55}{Ich kann mit dem Satz vom Nullprodukt eine Sachfrage beantworten.}
\end{abhakseite}

\end{document}
